% ------------------------------------------------------------------------------
% §8 Robustness — drafted Phase C from:
%   findings_3b.md (LOMO, June split); 03 findings (deviations, share row);
%   04 findings (treatments); 01 findings (thresholds, audits); rq1_wcb.csv
% ------------------------------------------------------------------------------
\section{Robustness}\label{sec:robustness}

\subsection{Leave-one-month-out}

Both headline unit-level RQ1 results were re-estimated 72 times, dropping one
month at a time under both outcome definitions. The Torrens coefficient is
negative in all 72 re-estimates; every flagged month is a $p$-value crossing
in the 0.03--0.15 band, never a sign change. Pelican Point fails the same
diagnostic: dropping October 2023 alone flips its sign (from $+0.184$ to
$-0.083$ fixed; $+0.128$ to $-0.105$ cost-indexed), and dropping two other
months moves it across the 5\% threshold in the opposite direction. The
asymmetry is the basis for the treatment in Section~\ref{sec:results}: Torrens
is a stable borderline-significant response; Pelican Point's positive
coefficient is a one-month artifact and is withdrawn.

\subsection{June 2022}

The crisis month is handled segment-by-segment using AEMO's own per-interval
flags (suspension window 15--24 June; administered pricing from 5 June). For
the RQ1 Torrens result: dropping the suspension window alone leaves the
coefficient intact ($-0.068$, $p=0.054$ cost-indexed), as does dropping the
pre-suspension fortnight alone ($-0.058$, $p=0.045$); only removing all of
June --- 19\% of Torrens's essential unit-rows, with zero essential rows after
the suspension --- degrades it to $p=0.13$. The dependence is on June's
essential-interval mass, a statistical-power statement, not on
suspension-era administered pricing, a confounding story. The RQ2 interaction
needs no such parsing: all four June treatments give $-0.044$ to $-0.058$
($p=0.001$--$0.029$; Figure~\ref{fig:rq2stability}).

\subsection{The timing wedge: separating the payment from the fuel}

The formula's memory creates a test the referee's fuel-confound story must
pass: gas prices fell in October 2022, but the trailing-percentile $\dt$
kept paying crisis rates until July 2023. If the essential-versus-matched
gap were fuel-stress conduct wearing the payment's clothes, it should have
collapsed with gas, nine months before the payment fell. It did not. With
the sample split at the calendar boundaries fixed in advance --- early 2022
($\dt\approx\$176$, a placebo-in-time), the crisis quarter (gas and $\dt$
both high), the lag-wedge (gas low, $\dt\approx\$340$), and the
post-roll-off period (both low, omitted) --- the eligibility-margin gap is
$-0.100$ in the crisis quarter ($p=0.03$), $-0.075$ through the lag-wedge
(bootstrap $p=0.089$), statistically indistinguishable from the crisis
value, and zero in both low-payment regimes despite their opposite fuel
conditions. The four period gaps order monotonically with the payment, not
with gas. What keeps this supportive rather than decisive is power: the
lag-wedge holds four essential-bearing months, and its contrast clears only
the 10\% level --- hence the ``consistent with'' language maintained
throughout this paper.

A conditional horse race points the same way within its limits. Adding
essential $\times$ gas alongside essential $\times$ $\dt$ (series
correlated at 0.48 across months), the payment interaction retains its full
magnitude and significance on the cheap-share outcome ($-0.054$, $p=0.013$)
with the gas interaction at exactly zero ($p=0.85$); on the eligibility
margin the two collinear series split the variance and neither is
individually significant --- a limitation of conditioning, which is why the
timing design above, not the horse race, carries the identification weight.

A placebo unit --- often essential but outside the compensation channel ---
does not exist in this market, for a reason that is itself informative:
essentiality concentrates on precisely the units the direction channel
pays. The candidate stations with constructed essentiality flags
(Quarantine, Dry Creek, Mintaro, Barker Inlet, Snapper Point) have 0--45
essential intervals in three years against Torrens's 4,083, and most are
themselves directed. The treatment set and the missing placebo have the
same cause.

\subsection{Inference and measurement}

Webb-weight bootstrap $p$-values match Rademacher throughout (RQ2 base case:
0.0040 versus 0.0041 fixed; 0.0061 versus 0.0073 cost-indexed). Because the
unrestricted wild bootstrap can be liberal with few unbalanced clusters,
the eligibility-margin estimate was additionally subjected to randomization
inference that imposes the null by construction --- permuting the
month-to-price assignment across the 36 sample months, 999 draws --- giving
$p=0.046$ (Table~\ref{tab:decomposition}); the null-imposed bootstrap
proper requires a compilation toolchain unavailable in the analysis
environment and remains a pre-submission upgrade. The two
outcome definitions agree on 96--100\% of interval classifications and on
every substantive conclusion. Threshold sweeps (\$150--\$500 fixed;
$1.5\times$--$3\times$ SRMC) leave the Torrens classification in an 82--96\%
band. The non-synchronous share specification (restricted to demand above
500~MW, avoiding the zero-crossing) gives $-0.017$/$-0.064$ for the pooled
essentiality coefficient --- same sign, same conclusion as the level
specification. The essentiality flag survives its leakage audit
($R^{2}\leq 0.0028$ on the unit's own behaviour) and its degeneracy audit
(essential intervals include units bidding as usual).

\subsection{An earlier design, reported not buried}

Before the matched interval design of Section~\ref{sec:design}, this project
ran a coarser opportunity-set version of the same question: does a withheld
indicator interact with $\dt$ on intervals where withholding could plausibly
cause a direction? That design is null (pooled interaction $t=1.55$,
$p=0.14$; one unit marginal at $p=0.07$), and Appendix~\ref{app:withhold_opportunity}
reports it in full. Two reasons it is superseded rather than averaged in: the
binary withheld outcome discards the intensive margin where
Section~\ref{sec:results}'s effect lives, and its opportunity flag conditions
on system state less sharply than the CEM match. The null is consistent with
the pattern in Section~\ref{sec:mechanism} --- the extensive-margin posture
does not flex with the prize; its depth does.
